\cotefile{157-5}
\scanpage{101}
\mspage{\unsure{88|98}}
% archive sheet number: 100

\uline{Dém} La nécessité de l'axiome de factorisation
pour $(TC,F)$, pour avoir un \unsure{triple} de Quillen,
est tautologique. Pour la suffisance, ayant
déjà les deux \unsure{principes|couples|quintuples} de Quillen $(C,TF),(TC,F)$
avec
\diagram{\begin{tikzcd}[ampersand replacement=\&,row sep=small,column sep=large]
TC \arrow[d,leftrightarrow] \&[-1em] \subset \&[-1em] C \arrow[d,leftrightarrow]\\
F \& \unsure{\supset|\longrightarrow} \& TF
\end{tikzcd}}
, et $TC=C\cap W$, le fait que \deleted{\unsure{c'est une condition}} \inserted{$(W,C,F)$ est}
un triple de Quillen est le cor.~4 p.~35
de $\overline{\underline{\mathrm{XV}}}$) : les conditions $TC=C\cap W$, $TF=F\cap W$,
$W=TF\circ TC$ sont équivalentes, et équivalent au
fait que $(W,C,F)$ soit un triple de Quillen.

Par la définition de Quillen des triples \uline{clos},
\unsure{par|vu} les propriétés
\[
\underbrace{W=TF\circ TC}_{\substack{\text{vraie pour tous}\\\text{les triples}\\\text{de Quillen}}},\qquad
\deleted{\unsure{E|F}}C=(TF)^{*},\qquad
\underbrace{F=(TC)_{\unsure{*|x}}}_{\substack{\text{vraie dans}\\\text{le cas actuel}}},
\]
il reste à exprimer que $C=(TF)^{*}$, ce qui
équivaut à : $C$ stable par facteurs directs\inserted{(cf loc. cit. cor.~1 p.~31)},
qed.

\uline{Remarques} 5.4. Les \deleted{\unsure{prop.}} énoncés 5.2. et 5.3,
et la relation $T\widetilde C\subset (TC)^{\unsure{4|\natural}}$ \unsure{aux|au} prop.~5.1,
ne font pas appel à la totalité de l'hypothèse que $(M,W,C)$ soit une catégorie à cofibrations spéciale, mais seulement que l'on ait
une cat. à cofibrations avec
\diagram{\begin{tikzcd}[ampersand replacement=\&,column sep=small,row sep=tiny]
C \arrow[r,leftrightarrow,"\shortmid" description] \& TF \arrow[d,equals,"\mathrm{déf}"] \&[-1em] \subset W\\
\& C_{*} \&
\end{tikzcd}}
\unsure{;|,}

